\chapter{Conclusion}
\label{ch:conclusion}

This report set out to explain the structure of the robot program, what each part
does, and how it developed, including the challenges we faced.

The program is split into five files, joined by \code{build.py} into the single
file the Brain can run. Autonomous movement uses measured geometry, \code{HOLD},
and proportional timeouts, which together fixed the original movement errors.
Turns are calibrated on the robot, and squaring to the wall gives the S--A--B
routine a repeatable start. Driver control shapes the stick input with squaring,
slew limiting and turn tapering, so the robot is smooth and predictable to drive.
The AI Vision sensor recognises the cone as a yellow--blue colour code, with a
fallback for when the code is missed.

The most important lessons came from the challenges. The vision sensor failed
without any error, and only on-screen diagnostics, each ruling out one cause, could
find the two settings that were wrong. The attempt to drive to the cone
autonomously showed that a sensor-driven routine is only as reliable as its
readings, and after 42 commits we chose a simpler, predictable routine instead.
However, the best parts of that work were kept as the driver's cone assist, where a
person can decide when to use it.

The main limitations are open-loop turns and the missing competition template. Both
have clear next steps, described in \cref{ch:evaluation}.
